% Ad Atticum 5.2 --- headnote
% ad-atticum-05-02, 10 May 51 BC, written from the Pompeian villa

\textit{Cicero to Atticus, written on 10 May 51 BC from the Pompeian
villa on the day of his departure southward toward Brundisium
and the unwanted Cilician command. The first letter in the
manuscript order of book 5 of the Atticus correspondence after
the missing-or-displaced opening; chronologically the second of
the journey-south series, written the morning he was leaving
Pompeii for Trebula. The body opens with the day's itinerary
and a report of a visit at the Cumanum from Q. Hortensius
Hortalus --- which Cicero treats as a piece of good fortune,
because Hortensius is offering himself as the man at Rome who
will press, when the time comes, that Cicero's provincial
command not be extended beyond the year the \textit{lex Pompeia}
allows. Atticus is asked to second the courtesy at Rome, and
the tribune-designate C. Furnius is being lined up on the same
cause.}

\textit{Section 2 is a small Pompeian comedy of manners: amid the
crowd of grandees at Cumae who came in numbers to see Cicero,
M. Vestorius's man Rufio managed not to call on him --- a
strategem (\textit{strategemate} is the Greek-derived word) for
which Cicero claims to give the man credit, since he made no
effort to be praised. The pun is dry; Cicero salutes him in
passing at the Puteoli market and lets the slight pass. \S 3
turns to the real consolation: he hopes the burden will last
no longer than a single year. Atticus is asked, when he is back
from Epirus, to put his weight behind keeping the appointment
short. The closing paragraph is the political weather report ---
Caesar's reaction to a recent decree, a worrying rumour about
the Transpadanes electing \textit{quattuorviri} on Roman
municipal lines --- with one obelized verb that the manuscript
has not yielded up; Pompey, who is at Tarentum, is to be the
source of better information. The launch-prompt summary
mentions ``Pomptina''; this is a confusion for Pomptinus, the
legate Cicero is waiting on at Brundisium, who is named in the
adjacent letters but not in 5.2 itself.}
